% Consolidated guide-ch12 from 9a1384ddf1afcd855b09256c9c6f1c90b5023140; source fragments preserved.

% === BEGIN guide-ch12-examples.tex blob 07f27d37c3d2bde832e3a3b75cd787fe76a7b6e1 ===
% 《大学物理学学习指导》第5版，第12章“光的干涉”例12.1—12.8。
% 书内第179—184页，源PDF物理第187—192页，按高分辨率图页核验。
% 不复制原书原文整段和受保护插图；全部推导原创中文整理。

\sourceexample{12.1}{179}{187}
\problemtext 杨氏双缝干涉中，分别改变屏幕距离、双缝间距、
波长，判断条纹疏密；其中一缝覆盖折射率$n>1$的透明片时，
零级明纹朝哪侧移动？
\workings
旁轴近似下条纹间隔
\[
\boxed{\Delta x=\frac{D\lambda}{d}}.
\]
因此屏距$D$减小则\emph{变密}；
缝距$d$减小则\emph{变疏}；
波长$\lambda$变长则\emph{变疏}。
将透明薄片厚度$t$放在下方缝$S_2$，
该路光程增加$(n-1)t$。
原中央明纹要重新满足等光程，屏上观察点必须更靠近被覆盖的$S_2$缝，
所以\textbf{零级与整个条纹系朝有薄片的缝一侧移动}。
\crosscheck 透明片使全部条纹整体平移而不改变（理想等厚均匀薄片下）
两缝几何条纹间隔；不要把条纹位移与间距变化混为一谈。

\sourceexample{12.2}{179--180}{187--188}
\problemtext 空气$n_1=1$、MgF$_2$膜$n_2=1.38$、
玻璃基底$n_3=1.50$构成减反结构；
对白光中的$\lambda=0.55\,\mu$m光，
求反射最弱时膜的最小正厚度。
\workings
由于$n_1<n_2<n_3$，
膜上、下界面两束主要反射光均发生$\pi$相变，
相对反射相位差不含额外半波项。
要使反射相消，膜内双程光程
$2n_2t=(k+\tfrac12)\lambda$。
取$k=0$，
\[
\boxed{t_{\min}=\frac{\lambda}{4n_2}
=\frac{550\,\mathrm{nm}}{4(1.38)}
\approx99.6\,\mathrm{nm}}.
\]
\crosscheck 相位满足只保证对该波长的反射\emph{减弱}，
而完全消光还需要两反射振幅匹配；
理想单层减反折射率为$\sqrt{n_1n_3}\simeq1.225$，
故本膜未必达到严格零反射。

\sourceexample{12.3}{180--181}{188--189}
\problemtext 两玻璃板间形成上表面向右升高的空气劈尖。
问上板向上平移、向右平移、增加楔角，
或把空气换成折射率$1.55$的均匀透明介质后，
反射等厚干涉条纹的移动及疏密变化。
\workings
小角劈尖$t(x)\simeq\theta x$，
固定反射干涉级次对应固定膜厚，条纹间隔
\[
\boxed{\Delta x\simeq\frac{\lambda}{2n\theta}}.
\]
(1) 上板上移令各点厚度增加，固定级次纹的位置向左（薄端）移动，
夹角未变则间距不变；增厚过多可能使有限视野中部分级次不再出现。
(2) 上板向右平移使固定膜厚位置向右移动，纹整体向右，间距不变。
(3) 楔角增大，固定级次位置向薄端移，间距减小，纹\textbf{变密}。
(4) 介质$n=1.55$，$\Delta x'=\Delta x/1.55$，纹\textbf{变密}。
\crosscheck 更换介质还可能改变反射界面处的相位损失，
使接触处暗/亮互换；“间距变密”不依赖这一相位取值。
图示移动方向按原书把劈尖薄端置于左边的约定。

\sourceexample{12.4}{181}{189}
\problemtext 杨氏双缝间距$d=0.4$ mm，屏到缝距离$D=0.5$m，
波长$\lambda=6.4\times10^{-5}$cm，
观察点离中央明纹$x=0.1$mm，求两束光的相位差。
\workings
$\lambda=6.4\times10^{-7}$m，
小角下几何光程差
\[
\Delta \simeq\frac{dx}{D}
=\frac{(4.0\times10^{-4})(1.0\times10^{-4})}{0.5}
=8.0\times10^{-8}\,\mathrm m.
\]
因两缝相干同相，
\[
\boxed{\Delta\phi=\frac{2\pi\Delta}{\lambda}
=\frac{2\pi(8.0\times10^{-8})}{6.4\times10^{-7}}
=\frac{\pi}{4}}.
\]
\crosscheck 相位差是角度量，不能只把光程差的数值当成相位差。

\sourceexample{12.5}{181--182}{189--190}
\problemtext 折射率$n_f=1.30$的油膜覆于折射率$n_s=1.50$的玻璃上，
垂直入射可调谐单色光，在波长500nm与700nm处都观测到
反射干涉相消，且其间无其他消光波长，求油膜厚度。
\workings
上表面空气$\to$油膜与下表面油膜$\to$玻璃均有半波相变，
相对相位因此没有额外半波项。
反射相消满足
$2n_ft=(k+\tfrac12)\lambda$。
设500nm与700nm对应整数$k_1>k_2$，
且两波长之间无其他消光，故$k_1=k_2+1$。
\[
( k_2+\tfrac32)500=(k_2+\tfrac12)700\quad(\mathrm{nm}).
\]
整理得$500k_2+750=700k_2+350$，
$k_2=2,k_1=3$，
于是
\[
\boxed{t=\frac{(3+\tfrac12)500}{2(1.30)}\,\mathrm{nm}
=\frac{1750}{2.6}\,\mathrm{nm}
\approx673.1\,\mathrm{nm}}.
\]
\crosscheck 代入700nm时$2nt=(2+\tfrac12)700=1750$nm，
两式一致；不能错把相邻级次取为相同的$k$。

\sourceexample{12.6}{182--183}{190--191}
\problemtext 三种透明介质组成的劈尖装置，图中上方折射率$1.52$，
劈尖膜折射率$1.62$；在接点$P$两侧，膜下材料分别为
$1.75$（左）及$1.52$（右）。
垂直入射单色光后，讨论接点两侧的反射等厚条纹
以及由接触点向外的亮暗序列。
\workings
在左侧，顶界面由$1.52\to1.62$反射发生$\pi$相变；
底界面由$1.62\to1.75$反射也发生$\pi$相变。
两束相对无额外半波损失，接触处$t=0$满足
$2nt=0$反射相长，故\textbf{左侧接点为亮纹}，
向外为亮、暗、亮、暗依次交替。
在右侧，顶界面反射仍有$\pi$相变；
底界面$1.62\to1.52$反射没有$\pi$相变，
两束产生一次相对半波差，$t=0$反射相消，
故\textbf{右侧接点为暗纹}，向外为暗、亮、暗、亮依次交替。
\crosscheck 该结论按原扫描图中各区域折射率
$1.52,1.62,1.75$及右侧$1.52$的相对排列得到；
若重绘图时互换区域，接点的亮暗也会交换。
在本条暂未复制版权原图，以文字标明区域和物理判断。

\sourceexample{12.7}{183}{191}
\problemtext 牛顿环两问：(1)波长600nm的第$k$条暗环
与450nm的第$k+1$条暗环重合，
透镜曲率半径190cm，求该暗环半径；
(2)波长500nm的第5条明环与未知波长的第6条明环重合，
求未知波长。
\workings
暗环满足$r_k^2=kR\lambda$，
故$k(600)=(k+1)(450)$，解得$k=3$。
\[
\boxed{r=\sqrt{3(1.90)(600\times10^{-9})}
\simeq1.85\times10^{-3}\,\mathrm m}.
\]
明环半径按第一亮环记$k=1$时
$r_k^2=(k-\tfrac12)R\lambda$，因此
\[
\boxed{\lambda_3
=500\,\mathrm{nm}\,\frac{5-\tfrac12}{6-\tfrac12}
\simeq409.1\,\mathrm{nm}}.
\]
\crosscheck 两问分别用暗环整数级次和亮环半整数级次，
不能把两种计数公式混用。

\sourceexample{12.8}{184}{192}
\problemtext 折射率$n=1.4$、楔角$\theta=10^{-4}$ rad的劈尖，
在垂直单色光照射下相邻明纹距$0.25$cm；
求真空波长。若劈尖长3.5cm，估算可见明条纹总数。
\workings
小楔角下$\Delta x=\lambda/(2n\theta)$，
\[
\boxed{\lambda=2n\theta\Delta x
=2(1.4)(10^{-4})(0.25\times10^{-2})
=7.0\times10^{-7}\,\mathrm m=700\,\mathrm{nm}}.
\]
用长$L=3.5$cm除以间距，
\[
\boxed{N\simeq\frac{L}{\Delta x}
=\frac{3.5}{0.25}=14}.
\]
\crosscheck 14代表在有限楔尖长度内约14个明纹周期。
若端点恰好为明纹且两端都计数，严格的\emph{条纹线条数}
还须说明端点取舍约定；原书采用“约14条”。
% === END guide-ch12-examples.tex ===
